\section{Preliminaries}\label{sec:preliminaries}

The Ekeland--Hofer--Zehnder capacity of a convex body is the least action
of a closed characteristic on its boundary. To compute it for polytopes,
we need a definition of characteristics at nonsmooth boundary points and
a way to compare curves without keeping them on the boundary. This section
introduces both: a Hamiltonian differential inclusion and Clarke's dual
action principle.

Throughout, a convex body is a compact convex subset of $\mathbb R^4$ with
nonempty interior. We place the origin in its interior when using its gauge
function or normalized facet equations. Translation does not change capacity.


\subsection{Symplectic conventions and action}\label{subsec:preliminaries-polytopes}

We write \(\mathbb{R}^4\) with coordinates \((q_1,q_2,p_1,p_2)\). The standard
complex structure is
\[
  J_0 =
  \begin{pmatrix}
    0 & -I_2 \\
    I_2 & 0
  \end{pmatrix},
\]
and the standard symplectic form is
\[
  \omega_0(u,v) = \langle J_0u,v\rangle
  = \sum_{i=1}^2 (u_{q_i}v_{p_i}-u_{p_i}v_{q_i}).
\]
We use the standard primitive
\[
  \lambda_0 = \frac{1}{2}\sum_{i=1}^2(q_i\,dp_i-p_i\,dq_i).
\]
Thus \(d\lambda_0=\omega_0\) and
\(\lambda_0|_x(v)=\frac12\langle -J_0v,x\rangle\).
For an absolutely continuous closed curve
\(\gamma\colon[0,T]\to\mathbb{R}^4\), its action is
\begin{equation}\label{eq:preliminaries-action}
  A(\gamma) \coloneqq \int_\gamma \lambda_0
  = \frac{1}{2}\int_0^T
    \langle -J_0\dot{\gamma}(t),\gamma(t)\rangle\,dt.
\end{equation}
This action is invariant under translation and orientation-preserving
reparametrization.
If \(\gamma\) bounds a smooth map \(u\colon D^2\to\mathbb{R}^4\), then Stokes'
theorem gives
\[
  A(\gamma)=\int_{D^2}u^*\omega_0.
\]
Thus action measures signed symplectic area.

Let \(K\subset\mathbb{R}^4\) be a smooth convex body with
\(0\in\operatorname{int}(K)\). At \(x\in\partial K\), the restriction of
\(\omega_0\) to the three-dimensional space \(T_x\partial K\) has a
one-dimensional kernel,
\[
  \ell_x
  =
  \{v\in T_x\partial K:\omega_0(v,w)=0
    \text{ for all }w\in T_x\partial K\}.
\]
This is the characteristic line. A \emph{closed characteristic} is a closed
curve on \(\partial K\) whose velocity lies in the positive characteristic
direction. We choose the positive direction by requiring
\(\lambda_0(\dot{\gamma})>0\).

The restriction \(\alpha=\lambda_0|_{\partial K}\) is a contact form. Its Reeb
vector field \(R\) is the contact-normalized characteristic direction,
\[
  \alpha(R_x)=1,
  \qquad
  \omega_0(R_x,v)=0
  \quad\text{for all }v\in T_x\partial K.
\]
A Reeb orbit is a periodic integral curve of \(R\). It is the same geometric
object as a closed characteristic, with the parametrization fixed by
\(\lambda_0(\dot{\gamma})=1\). Consequently, if \(\gamma\) is a Reeb orbit on
\([0,T]\), then
\[
  A(\gamma)=\int_0^T\lambda_0(\dot{\gamma}(t))\,dt=T.
\]


\subsection{Convex Hamiltonians}
\label{subsec:preliminaries-convex-hamiltonian-language}

The smooth characteristic line field has a convex-analytic formulation that
continues to make sense when \(\partial K\) is not smooth.  For a convex body
\(K\) with the origin in its interior, its support and gauge functions are
\[
  h_K(y)\coloneqq\max_{x\in K}\langle x,y\rangle,
  \qquad
  g_K(x)\coloneqq\inf\{r>0:x\in rK\}.
\]
The support function gives the supporting height in direction $y$;
the gauge measures the dilation of $K$ needed to reach $x$. Thus the boundary
is the level set \(g_K=1\). We use the convex,
\(2\)-homogeneous Hamiltonian
\[
  H_K(x)\coloneqq g_K(x)^2.
\]
For a convex function \(H\), its subdifferential at \(x\) is
\[
  \partial H(x)
  =\{y:H(z)\geq H(x)+\langle y,z-x\rangle
    \text{ for every }z\}.
\]
When \(H\) is differentiable, this set consists only of \(\nabla H(x)\).

Here $W^{1,2}$ denotes the space of absolutely continuous curves with
square-integrable derivative.

\begin{definition}[Contact-normalized generalized characteristic]
\label{def:preliminaries-generalized-characteristic}
A \emph{contact-normalized generalized characteristic} on \(\partial K\) is
a curve \(\gamma\in W^{1,2}([0,T],\mathbb R^4)\), with \(T>0\) and
\(\gamma(0)=\gamma(T)\), whose image lies in \(\partial K\) and which satisfies
\begin{equation}\label{eq:preliminaries-generalized-characteristic}
  -J_0\dot\gamma(t)\in\partial H_K(\gamma(t))
  \quad\text{for almost every }t\in[0,T].
\end{equation}
\end{definition}

For smooth \(K\), the subdifferential in
\eqref{eq:preliminaries-generalized-characteristic} is the singleton
\(\{\nabla H_K\}\).  The resulting Hamiltonian vector is tangent to the
characteristic line on \(\partial K\).  Thus the definition recovers the Reeb
orbits described above, including their parametrization.

The normalization also persists without smoothness. Indeed, if
\(y\in\partial H_K(x)\), then the \(2\)-homogeneity of \(H_K\) gives
\[
  \langle y,x\rangle=2H_K(x).
\]
To obtain this identity, apply the subgradient inequality at $x$ to $rx$,
use $H_K(rx)=r^2H_K(x)$, and let $r$ approach $1$ from either side.
Along a generalized characteristic, \(H_K(\gamma)=1\), and hence
\[
  \lambda_0|_{\gamma(t)}(\dot\gamma(t))
  =\frac12\langle-J_0\dot\gamma(t),\gamma(t)\rangle
  =1
  \quad\text{for almost every }t.
\]
Consequently every contact-normalized generalized characteristic satisfies
\begin{equation}\label{eq:preliminaries-generalized-characteristic-action-period}
  A(\gamma)=T.
\end{equation}

For a polytope, we will compute $\partial H_K$ directly from its facet
normals. At a facet interior it is a singleton; at an intersection of facets
it is their convex hull, with the normalization above.

\subsection{EHZ capacity}\label{subsec:preliminaries-smooth-symplectic-geometry}

For smooth convex bodies, the first Ekeland--Hofer capacity and the
Hofer--Zehnder capacity
coincide. Their common value is the minimum action among closed
characteristics on the boundary. We use the resulting EHZ capacity in
the form
\begin{equation}\label{eq:preliminaries-ehz-capacity}
  c_{\mathrm{EHZ}}(K)
  = \min\{A(\gamma):\gamma
    \text{ is a closed characteristic on }\partial K\}.
\end{equation}
This is the smooth least-action characterization of
\cite[Theorem~2.2]{AAO2014}.
The minimum is positive and attained. By the preceding normalization,
it is also the least period of a Reeb orbit.

For a nonsmooth convex body, we define capacity as the limit of the
capacities of smooth convex approximations in the Hausdorff metric. This
limit is independent of the approximating sequence. The least-action
description remains valid when closed characteristics are replaced by the
generalized characteristics of
Definition~\ref{def:preliminaries-generalized-characteristic}
\cite[Proposition~2.7]{AAO2014}. Their finite
facet description on a polytope is derived in
Section~\ref{sec:generalized-reeb-orbits-polytopes}.

As a symplectic capacity, \(c_{\mathrm{EHZ}}\) is monotone
under inclusion, invariant under symplectomorphisms, and conformal:
\[
  c_{\mathrm{EHZ}}(\lambda K)=\lambda^2c_{\mathrm{EHZ}}(K)
  \quad(\lambda>0).
\]
It is normalized by
\[
  c_{\mathrm{EHZ}}(B^4(r))=c_{\mathrm{EHZ}}(Z^4(r))=\pi r^2,
\]
where \(B^4(r)\) is the radius-\(r\) ball and
\(Z^4(r)=B^2(r)\times\mathbb{R}^2\) is the symplectic cylinder, whose
disc factor occupies a symplectic coordinate plane.

Throughout the thesis, \(\operatorname{vol}_4\) denotes four-dimensional
Lebesgue measure in the coordinate order \((q_1,q_2,p_1,p_2)\). In dimension
four we define the systolic ratio by
\begin{equation}\label{eq:preliminaries-systolic-ratio}
  \operatorname{sys}(K)
  \coloneqq
  \frac{c_{\mathrm{EHZ}}(K)^2}{2\operatorname{vol}_4(K)}.
\end{equation}
With this normalization, Viterbo's proposed inequality for four-dimensional convex
bodies was
\[
  \operatorname{sys}(K)\leq 1,
  \qquad\text{equivalently}\qquad
  c_{\mathrm{EHZ}}(K)^2\leq 2\operatorname{vol}_4(K)
\]
\cite{Viterbo2000}.

\begin{definition}[Systolic-ratio symmetries]
\label{def:preliminaries-systolic-ratio-symmetries}
Let
\[
  G_{\operatorname{sys}}
  =
  \{x\mapsto b+\rho Mx:
    b\in\mathbb R^4,\ \rho>0,\ M\in\operatorname{Sp}(4,\mathbb R)\}.
\]
This is the group of affine transformations generated by translations,
positive scalings, and linear symplectic transformations.
\end{definition}

\begin{lemma}[Systolic-ratio invariance]
\label{lem:preliminaries-systolic-ratio-invariance}
For every convex body \(K\subset\mathbb R^4\) and every
\(g\in G_{\operatorname{sys}}\),
\[
  \operatorname{sys}(gK)=\operatorname{sys}(K).
\]
\end{lemma}

\begin{proof}
Translations and linear symplectic transformations are symplectomorphisms, so
they preserve \(c_{\mathrm{EHZ}}\). They also preserve four-dimensional
volume, because translations do not change volume and every symplectic matrix
has determinant \(1\). A positive scaling by \(\rho\) gives
\[
  c_{\mathrm{EHZ}}(\rho K)=\rho^2c_{\mathrm{EHZ}}(K),
  \qquad
  \operatorname{vol}_4(\rho K)=\rho^4\operatorname{vol}_4(K).
\]
Substituting these identities in
\eqref{eq:preliminaries-systolic-ratio} shows that the ratio is unchanged.
\end{proof}

\begin{proposition}[Continuity]\label{prop:preliminaries-sys-hausdorff-continuity}
Capacity and systolic ratio are continuous on the space of convex bodies
with respect to the Hausdorff metric.
\end{proposition}

\begin{proof}
Fix a convex body $K$ and translate it so that $rB^4(1)\subset K$ for
some $r>0$. If $K'$ has Hausdorff distance at most $\delta<r$ from $K$,
its support function differs from $h_K$ by at most $\delta$ on the unit
sphere. Since $h_K\geq r$ there, this gives
\[
 (1-\delta/r)K\subset K'\subset(1+\delta/r)K.
\]
Monotonicity and conformality squeeze $c_{\rm EHZ}(K')$ between
$(1-\delta/r)^2c_{\rm EHZ}(K)$ and
$(1+\delta/r)^2c_{\rm EHZ}(K)$. The same inclusions squeeze volume between
the corresponding fourth powers times $\operatorname{vol}_4(K)$.
Both quantities therefore converge to their values at $K$. Since volume
is positive, their quotient defining the systolic ratio is continuous too.
\end{proof}

\subsection{Facets and polar vertices}
\label{subsec:preliminaries-polytope-input-language}

For a polytope, we absorb each supporting height into its normal vector. Let
\(K\subset\mathbb R^4\) be full-dimensional with
\(0\in\operatorname{int}(K)\). We write
\begin{equation}\label{eq:preliminaries-normalized-halfspaces}
  K=\{x\in\mathbb R^4:\langle a_i,x\rangle\leq1,
    \ i=1,\ldots,F\}.
\end{equation}
If \(n_i\) is the outward unit normal to the corresponding supporting
hyperplane and \(h_i=h_K(n_i)>0\) is its height, then
\[
  a_i=\frac{n_i}{h_i}.
\]
For an irredundant description, the vectors \(a_i\) are precisely the vertices
of the polar body
\[
  K^\circ=\{y:\langle y,x\rangle\leq1\text{ for all }x\in K\}
  =\operatorname{conv}\{a_1,\ldots,a_F\}.
\]


Conversely, the normalized halfspaces define a bounded full-dimensional
polytope containing the origin in its interior exactly when
\(0\in\operatorname{int}\operatorname{conv}\{a_i\}\).  The description is
irredundant exactly when the listed vectors are pairwise distinct and every
\(a_i\) is an extreme point of this convex hull.

The ordered list \((a_1,\ldots,a_F)\) is a labelled presentation.  Relabelling
the same inequalities does not change the unlabelled convex body, but labels
matter when an algorithm records facet words or when a local calculation uses
the \(a_i\) as coordinates.

\subsection{Lagrangian products}
\label{subsec:preliminaries-lagrangian-products}

Write
\[
  \mathbb R^4=\mathbb R_q^2\oplus\mathbb R_p^2
  =\{(q_1,q_2,p_1,p_2)\}.
\]
\begin{definition}[Lagrangian product]
\label{def:preliminaries-lagrangian-product}
For planar convex bodies \(K_q,K_p\subset\mathbb R^2\), their Lagrangian
product is
\begin{equation}\label{eq:preliminaries-lagrangian-product}
  K_q\times_L K_p
  =\{(q_1,q_2,p_1,p_2):(q_1,q_2)\in K_q,\ (p_1,p_2)\in K_p\}.
\end{equation}
Thus \(K_q\) occupies the Lagrangian \(q\)-plane and \(K_p\) occupies the
Lagrangian \(p\)-plane.  This is different from taking a product of planar
domains in the symplectic coordinate planes \((q_1,p_1)\) and \((q_2,p_2)\).
\end{definition}

The symplectic form vanishes on each factor plane separately, and pairs
the two factors by
\[
 \omega_0((u,0),(0,v))=\langle u,v\rangle.
\]
Consequently a Hamiltonian normal in one factor generates motion in the
other: $J_0(u,0)=(0,u)$ and $J_0(0,v)=(-v,0)$. This coupling is central
to the capacity formulas for products of polygons.

\subsection{Clarke's dual action principle}
\label{subsec:preliminaries-clarke-dual-action-principle}

A curve constrained to lie on $\partial K$ is difficult to modify: changing
one velocity can move the rest of it off the boundary. We will instead work
with closed curves in $\mathbb R^4$ and a functional that depends only on
velocity. The dual action principle says that minimizing this functional
recovers the capacity, and that a minimizer can be placed back on the
boundary. This is the freedom used in the next section to rearrange facet
velocities.

Recall $H_K=g_K^2$, and define its convex conjugate by
$H_K^*(y)=\sup_x(\langle x,y\rangle-H_K(x))$. For a closed curve
$z\in W^{1,2}([0,T],\mathbb R^4)$ put
\begin{equation}\label{eq:preliminaries-dual-functional}
 I_K(z)=\int_0^T H_K^*(-J_0\dot z(t))\,dt
       =\frac14\int_0^T h_K(-J_0\dot z(t))^2\,dt.
\end{equation}
The conjugate identity will be proved below. We impose the action constraint
$A(z)=T$ and minimize over all periods:
\begin{equation}\label{eq:preliminaries-dual-problem}
 d(K)=\inf\{I_K(z):T>0,\ z(0)=z(T),\ A(z)=T\}.
\end{equation}
The infimum is over the closed Sobolev curves just defined. A characteristic
with its contact parametrization is feasible because its action equals its
period. For a general curve, the constraint fixes its spatial scale relative
to its period; its image need not lie on $\partial K$.

\begin{theorem}[Dual action principle]
\label{thm:preliminaries-clarke-dual-action}
For a convex body $K\subset\mathbb R^4$ with $0\in\operatorname{int}K$,
the infimum in \eqref{eq:preliminaries-dual-problem} is attained and
\[
 d(K)=c_{\rm EHZ}(K)
     =\min_{\gamma\subset\partial K}A(\gamma),
\]
where the last minimum is over contact-normalized generalized
characteristics. Every such characteristic satisfies $I_K(\gamma)=A(\gamma)$.
Every dual minimizer can be rescaled and translated into a minimizing
generalized characteristic, without changing its value of $I_K$.
If a dual minimizer already has $I_K(z)=T$, only a translation is needed.
\end{theorem}

We use the nonsmooth variational results underlying Clarke's principle
\cite[Proposition~2.5 and Lemmas~5.1--5.2]{AAO2014}; for polytopes the
correspondence is also stated in \cite[Lemma~2.1]{HK2017}. The following
argument explains the conjugacy, normalization and reconstruction used here.

\paragraph{The conjugate and its equality condition.}
\begin{lemma}\label{lem:preliminaries-fenchel-duality}
One has $H_K^*(y)=h_K(y)^2/4$, and
\[
 H_K(x)+H_K^*(y)\geq\langle x,y\rangle.
\]
Equality holds exactly when $y\in\partial H_K(x)$, equivalently
$x\in\partial H_K^*(y)$.
\end{lemma}
\begin{proof}
Write $x=ru$ with $r\geq0$ and $g_K(u)=1$. For a fixed $r$, the largest
value of $\langle x,y\rangle$ is $r h_K(y)$. Thus
\[
 H_K^*(y)=\sup_{r\geq0}(r h_K(y)-r^2)=\tfrac14h_K(y)^2.
\]
The inequality follows from the definition of the supremum. Equality says
that $x$ maximizes $\langle x,y\rangle-H_K(x)$, which is precisely the
subgradient inequality defining $y\in\partial H_K(x)$. Applying conjugacy
in the other direction gives the equivalent condition.
\end{proof}

\paragraph{The variational equation.}
A minimizer satisfies a weak Euler--Lagrange equation. Its useful form is
\begin{equation}\label{eq:preliminaries-dual-euler-lagrange}
 \nu z(t)+b\in\partial H_K^*(-J_0\dot z(t))
 \quad\text{for almost every }t,
\end{equation}
for a constant vector $b$ and a real multiplier $\nu$.
\begin{definition}[Weak dual critical curve]
\label{def:preliminaries-weak-dual-critical}
A feasible curve
satisfying this equation is called \emph{weak dual critical}.
\end{definition}
The constant $b$ reflects the freedom to translate the curve. Conjugacy
rewrites the same equation as
$-J_0\dot z\in\partial H_K(\nu z+b)$, bringing it close to the
characteristic equation.

To see the form of the variational equation, choose
$p(t)\in\partial H_K^*(-J_0\dot z(t))$. The derivative of action in a
periodic variation $v$ is $DA(z)[v]=\int\langle-J_0\dot z,v\rangle$.
The cited nonsmooth multiplier theorem provides a measurable $L^2$ choice
of $p$ such that, for every smooth periodic variation $v$,
\[
 \int_0^T\langle p,-J_0\dot v\rangle\,dt
 =\nu\int_0^T\langle-J_0\dot z,v\rangle\,dt.
\]
The left side pairs the chosen subgradient with the change in velocity;
the right side is the variation of the action constraint.
Integration by parts gives
$-J_0\dot p=-\nu J_0\dot z$ in distributions, hence
$\dot p=\nu\dot z$ and $p=\nu z+b$. The action constraint is regular
because $DA(z)[z]=2T\ne0$. Thus the multiplier equation has the stated signs and constants.

For any weak critical curve, the multiplier is determined by its value.
The two-homogeneity of $H_K^*$ gives
$\langle p,y\rangle=2H_K^*(y)$ when $p\in\partial H_K^*(y)$.
Pairing \eqref{eq:preliminaries-dual-euler-lagrange} with $-J_0\dot z$
and integrating, the constant vector contributes zero because the curve
is closed. We obtain
\begin{equation}\label{eq:preliminaries-dual-multiplier-value}
 2I_K(z)=2\nu A(z)=2\nu T,
 \qquad \nu=\frac{I_K(z)}{T}>0.
\end{equation}
Positivity follows since a feasible curve has positive action and is
nonconstant, while $H_K^*$ is positive away from zero.

\paragraph{Reconstructing a boundary curve.}
Let $z$ be weak critical and set $x=\nu z+b$ and $y=-J_0\dot z$.
Then $y\in\partial H_K(x)$. The convex chain rule gives
\[
 \frac{d}{dt}H_K(x(t))=\langle y(t),\dot x(t)\rangle
 =\nu\langle-J_0\dot z(t),\dot z(t)\rangle=0.
\]
Thus $H_K(x)$ is constant. Since $y\in\partial H_K(x)$, Fenchel equality gives
$H_K(x)+H_K^*(y)=\langle x,y\rangle$. Two-homogeneity gives
$\langle x,y\rangle=2H_K(x)$, so $H_K^*(y)=H_K(x)$ almost everywhere.
Integrating the constant value yields
\[
 I_K(z)=\int_0^T H_K^*(y)\,dt=T H_K(x),
 \qquad H_K(x)=I_K(z)/T=\nu.
\]
Put $c=\sqrt\nu$ and define
\[
 \gamma(t)=c\,z(t/c^2)+b/c,\qquad 0\leq t\leq c^2T.
\]
Since $\gamma=x(t/c^2)/c$, it satisfies $H_K(\gamma)=1$.
Subgradient homogeneity also gives
\[
 -J_0\dot\gamma(t)=y(t/c^2)/c\in\partial H_K(\gamma(t)).
\]
This is the contact-normalized characteristic equation. Its period and
action are $c^2T=I_K(z)$. Direct substitution in
\eqref{eq:preliminaries-dual-functional} gives $I_K(\gamma)=I_K(z)$.
In particular, when $I_K(z)=T$ we have $c=1$, so reconstruction only
translates the curve.

Conversely, a contact-normalized characteristic satisfies
$H_K(\gamma)=1$ and $-J_0\dot\gamma\in\partial H_K(\gamma)$.
Fenchel equality therefore gives $H_K^*(-J_0\dot\gamma)=1$ almost
everywhere. Consequently $I_K(\gamma)=T=A(\gamma)$, and it is a feasible
dual curve. An attained dual minimum reconstructs to a characteristic
with the same value; every characteristic supplies a feasible dual value.
This proves equality of the minima and, by the least-action
characterization, their equality with capacity.

\paragraph{Period convention and attainment.}
Haim--Kislev formulate the principle using $w:[0,1]\to\mathbb R^4$ with
$A(w)=1/2$. If $z$ is feasible in our convention, then
\[
 w(s)=\frac{z(Ts)}{\sqrt{2T}},\qquad I_K(w)=\tfrac12 I_K(z).
\]
Conversely, $z(t)=\sqrt{2T}\,w(t/T)$ is feasible for any chosen $T>0$
and has $I_K(z)=2I_K(w)$. If the fixed-period formulation imposes mean zero,
subtract the mean of $w$; translation changes neither its action nor $I_K$.
Thus the fixed-period attainment theorem applies
to our infimum, with its factor of two accounted for. Choosing
$T=2I_K(w)$ gives an attained representative with $I_K(z)=T$.

The distinction needed below is now concrete. Feasibility allows us to
compare functional values, but does not place a curve on the boundary.
Once a modified curve is shown to minimize $I_K$, the variational equation
and reconstruction justify placing it there. This is why rearranging
velocities will be done in the dual problem first.
